\section{Baseline và quy ước thiết kế Phase 2}

\subsection{Mục tiêu và phạm vi}

Phase 2 chuyển baseline yêu cầu của Submission 1 thành thiết kế giao diện và mô hình hành vi có thể kiểm tra. Báo cáo gồm UI mockup ở mức nhóm, sequence diagram và activity diagram cho 18 use case; state-chart diagram được trình bày riêng như phần bonus.

Phase 2 không tự ý thêm actor, business rule, trạng thái hay nghiệp vụ mới. Nếu phát hiện thay đổi cần thiết, nhóm phải cập nhật baseline Phase 1 và truy vết lại các tài liệu liên quan trước khi xem thay đổi là hợp lệ.

\subsection{Nguồn truy vết từ Submission 1}

\begin{table}[htbp]
  \centering
  \caption{Nguồn truy vết bắt buộc của thiết kế Phase 2}
  \begin{tabularx}{\textwidth}{|L{34mm}|Y|Y|}
    \hline
    \rowcolor{gray!15}
    \textbf{Nguồn Phase 1} & \textbf{Cách sử dụng trong Phase 2} & \textbf{Dấu vết trong báo cáo}\\\hline
    Actor Catalog & Xác định vai trò thao tác và hệ thống ngoài tham gia. & Actor/lifeline/swimlane giữ nguyên mã \texttt{ACT-*}.\\\hline
    Use Case Specification & Cung cấp trigger, precondition, main flow, alternative flow và exception flow. & Mỗi cặp \texttt{SEQ-UC-*}, \texttt{ACT-UC-*} trỏ về một \texttt{UC-*}.\\\hline
    Business Rule và Requirement & Xác định guard, validation, tác vụ tự động và phản hồi lỗi. & Ghi rõ \texttt{BR-*}, \texttt{FR-*}, \texttt{NIFR-*}, \texttt{NFR-*} liên quan.\\\hline
    Thuật ngữ và trạng thái & Giữ nhất quán tên đối tượng và trạng thái giữa UI, diagram và code sau này. & Không tự thay \texttt{Available}, \texttt{Confirmed}, \texttt{OutOfService}, v.v. bằng từ đồng nghĩa.\\\hline
  \end{tabularx}
\end{table}

\subsection{Quy ước mã và tên}

\begin{itemize}
  \item Màn hình dùng mã \texttt{SCR-[VAI TRÒ]-[TÊN]}, ví dụ \texttt{SCR-STU-HUB-DETAIL}.
  \item Sequence diagram dùng mã \texttt{SEQ-UC-[PHÂN HỆ]-[SỐ]}.
  \item Activity diagram dùng mã \texttt{ACT-UC-[PHÂN HỆ]-[SỐ]}.
  \item State chart dùng mã \texttt{STM-[ĐỐI TƯỢNG]}.
  \item Tên message trong sequence diagram là hành động có nghĩa nghiệp vụ; không dùng tên nút giao diện làm tên service nếu chưa có quyết định thiết kế.
\end{itemize}

\subsection{Quy tắc mô hình hóa}

Sequence diagram nhấn mạnh thứ tự trao đổi theo thời gian; activity diagram nhấn mạnh quy trình, quyết định và trách nhiệm. Hai sơ đồ không thay thế nhau và phải cùng phản ánh use-case flow. Các nhánh \texttt{alt}/\texttt{opt} trong sequence diagram và guard trong activity diagram phải dùng điều kiện có thể đánh giá, không ghi chung chung như ``hợp lệ'' nếu baseline đã có business rule cụ thể.
